\documentclass[11pt,letterpaper]{article}

% ==============================================================================
% PAQUETES DE CONFIGURACIÓN Y TIPOGRAFÍA
% ==============================================================================
\usepackage[utf8]{inputenc}
\usepackage[spanish,es-nodecimaldot,es-noshorthands]{babel}
\usepackage[margin=2.2cm,top=2.5cm,bottom=2.5cm]{geometry}
\usepackage{amsmath,amssymb,amsfonts,mathtools,bm}
\usepackage{microtype}
\usepackage{booktabs,array,tabularx}
\usepackage{fancyhdr}
\usepackage{lastpage}
\usepackage{xcolor}
\usepackage{tikz}
\usetikzlibrary{arrows.meta,patterns,calc,positioning,decorations.pathmorphing,shapes.geometric,babel}
\usepackage[skins,breakable]{tcolorbox}
\usepackage{hyperref}

% ==============================================================================
% PALETA DE COLORES INSTITUCIONAL Y PEDAGÓGICA (TEMA ESMERALDA / UIS)
% ==============================================================================
\definecolor{uisgreen}{RGB}{5, 150, 105}       % Color principal Clase 04
\definecolor{uisgreenlight}{RGB}{236, 253, 245}
\definecolor{uisblue}{RGB}{2, 132, 199}
\definecolor{uispurple}{RGB}{109, 40, 217}
\definecolor{uisamber}{RGB}{217, 119, 6}
\definecolor{uisrose}{RGB}{220, 38, 38}
\definecolor{uisdark}{RGB}{15, 23, 42}
\definecolor{uisgray}{RGB}{100, 116, 139}
\definecolor{uislightbg}{RGB}{248, 250, 252}
\definecolor{uisboxborder}{RGB}{203, 213, 225}

% ==============================================================================
% HIPERVÍNCULOS
% ==============================================================================
\hypersetup{
    colorlinks=true,
    linkcolor=uisgreen,
    citecolor=uisblue,
    urlcolor=uisgreen,
    pdftitle={La Rotula Plastica ASCE 41 en Columna en Voladizo - Teoria y Ejemplo Numerico},
    pdfauthor={Prof. Orlando Arroyo - Universidad Industrial de Santander}
}

% ==============================================================================
% ENCABEZADOS Y PIES DE PÁGINA (FANCYHDR)
% ==============================================================================
\pagestyle{fancy}
\fancyhf{}
\renewcommand{\headrulewidth}{0.6pt}
\renewcommand{\footrulewidth}{0.4pt}
\fancyhead[L]{\small\textbf{\color{uisgreen}Análisis No Lineal de Estructuras} | Posgrado en Ing. Estructural}
\fancyhead[R]{\small\color{uisgray}Universidad Industrial de Santander}
\fancyfoot[L]{\small\color{uisgray}Prof. Orlando Arroyo}
\fancyfoot[C]{\small\color{uisgray}Material Docente de Cátedra --- Clase 04}
\fancyfoot[R]{\small\color{uisdark}\thepage\ / \pageref{LastPage}}

% ==============================================================================
% CAJAS DESTACADAS (TCOLORBOX)
% ==============================================================================
\newtcolorbox{conceptbox}[1][]{
    enhanced,
    breakable,
    colback=uislightbg,
    colframe=uisgreen,
    coltitle=white,
    fonttitle=\bfseries\sffamily,
    title={#1},
    sharp corners=downhill,
    arc=4mm,
    boxrule=1.2pt,
    left=12pt,right=12pt,top=8pt,bottom=8pt,
    shadow={2pt}{-2pt}{0pt}{black!10}
}

\newtcolorbox{stepbox}[1][]{
    enhanced,
    breakable,
    colback=white,
    colframe=uisblue,
    coltitle=white,
    fonttitle=\bfseries\sffamily,
    title={#1},
    rounded corners,
    arc=3mm,
    boxrule=1.2pt,
    left=12pt,right=12pt,top=8pt,bottom=8pt
}

\newtcolorbox{alertbox}[1][]{
    enhanced,
    breakable,
    colback=uisrose!5,
    colframe=uisrose,
    coltitle=white,
    fonttitle=\bfseries\sffamily,
    title={#1},
    rounded corners,
    arc=3mm,
    boxrule=1.2pt,
    left=12pt,right=12pt,top=8pt,bottom=8pt
}

\newtcolorbox{examplebox}[1][]{
    enhanced,
    breakable,
    colback=uisamber!5,
    colframe=uisamber,
    coltitle=white,
    fonttitle=\bfseries\sffamily,
    title={#1},
    rounded corners,
    arc=3mm,
    boxrule=1.2pt,
    left=12pt,right=12pt,top=8pt,bottom=8pt
}

% ==============================================================================
% DOCUMENTO PRINCIPAL
% ==============================================================================
\begin{document}

% --- PORTADA / ENCABEZADO TITULAR ---
\begin{center}
    {\scshape\Large Universidad Industrial de Santander}\\[0.2cm]
    {\large Escuela de Ingeniería Civil --- Posgrado en Ingeniería Estructural}\\[0.35cm]
    {\color{uisgreen}\hrule height 2pt}\vspace{0.4cm}
    {\huge\bfseries\color{uisdark} Mecánica de la Rótula Plástica Concentrada (ASCE 41)}\\[0.25cm]
    {\Large\bfseries\color{uisgreen} Acoplamiento en Serie, Control por Desplazamiento y Fenómeno de Descarga Elástica (\textit{Springback}) en Columna en Voladizo}\\[0.35cm]
    {\color{uisgreen}\hrule height 1pt}\vspace{0.35cm}
    {\textbf{Autor:} Prof. Orlando Arroyo}\quad\textbullet\quad
    {\textbf{Asignatura:} Análisis No Lineal de Estructuras}\quad\textbullet\quad
    {\textbf{Versión:} Documentación Técnica Canónica (2026)}
\end{center}

\vspace{0.25cm}

\begin{abstract}
\noindent El presente documento desarrolla de manera rigurosa los fundamentos físicos, cinemáticos y constitutivos de una \textbf{rótula plástica concentrada} según los lineamientos de las normas \textbf{ASCE/SEI 41-17} y \textbf{ASCE/SEI 41-23}, aplicada al caso fundamental de una columna en voladizo (\textit{cantilever}) de sección rectangular ($b \times d$) y longitud $L$. Se modela la disipación inelástica mediante un resorte rotacional no lineal ubicado a una distancia representativa del $5\%$ de la altura ($x_h = 0.05 L$), acoplado en serie con el fuste elástico continuo del elemento. Se deduce formalmente la relación de flexibilidad combinada y se explica en detalle el algoritmo matemático para determinar la fuerza lateral cortante $V$ correspondiente a un desplazamiento objetivo impuesto $\Delta$ en el marco de un análisis estático no lineal (\textit{Pushover}). Asimismo, se demuestra matemática y físicamente el fenómeno de \textbf{descarga elástica} (\textit{springback}) del fuste durante la etapa de ablandamiento post-pico (\textit{softening}) de la rótula. Finalmente, se presentan cuatro pasos de cálculo numérico manual paso a paso con verificación exacta del equilibrio seccional.
\end{abstract}

\vspace{0.35cm}

% ==============================================================================
% SECCIÓN 1: INTRODUCCIÓN Y MODELO DE PLASTICIDAD CONCENTRADA
% ==============================================================================
\section{Introducción y Modelo de Plasticidad Concentrada}

En el análisis inelástico de pórticos y elementos estructurales sometidos a acciones sísmicas severas, la disipación de energía por flexión se concentra en regiones finitas denominadas \textbf{zonas de rótula plástica}. Existen dos grandes enfoques computacionales para capturar este fenómeno:
\begin{enumerate}
    \item \textbf{Plasticidad Distribuida (\textit{Distributed Plasticity / Fiber Elements}):} El elemento se discretiza longitudinalmente en puntos de integración transversalmente divididos en fibras unidimensionales con relaciones constitutivas $\sigma - \varepsilon$ no lineales. Aunque es riguroso, exige un costo computacional elevado y es propenso a problemas de localización de deformaciones dependientes de la malla si no se aplican técnicas de regularización energética.
    \item \textbf{Plasticidad Concentrada (\textit{Lumped / Concentrated Plasticity}):} Toda la deformación inelástica se idealiza en un \textbf{resorte rotacional de longitud cero} con una ley momento-rotación ($M_h - \theta_h$) precalibrada, mientras que el resto del elemento se mantiene con comportamiento elástico lineal. Este enfoque es la base fundamental de la evaluación basada en desempeño estipulada en los estándares ASCE 41.
\end{enumerate}

\begin{conceptbox}[¿Por qué ubicar la rótula plástica a $x_h = 0.05 L$? Fundamento Físico]
En una columna en voladizo empotrada en la base y cargada lateralmente en la punta superior, el diagrama de momentos flectores es máximo en la base ($x=0$) y decrece linealmente hacia la punta ($x=L$). La zona de plastificación real del hormigón y fluencia del acero se extiende sobre una longitud finita $L_p$, comúnmente estimada mediante la fórmula semi-empírica de Paulay y Priestley (1992):
\begin{equation}
L_p \approx 0.08 L + 0.022 d_b f_y \approx 0.40 d \text{ a } 0.50 d \approx 0.10 L
\end{equation}
donde $d$ es el peralte de la sección y $d_b$ el diámetro del refuerzo longitudinal. La curvatura plástica real $\phi_p(x)$ se distribuye sobre este tramo. La cinemática equivalente de un resorte rotacional concentrado exige que la rotación plástica total $\theta_p = \int_0^{L_p} \phi_p(x) dx$ actúe en el \textbf{centroide del diagrama de curvatura inelástica}, el cual se sitúa aproximadamente en el punto medio de la zona plastificada:
\begin{equation}
x_h = \frac{L_p}{2} \approx 0.05 L
\end{equation}
Ubicar el resorte rotacional exactamente al $5\%$ de la altura captura con alta fidelidad la rotación de cuerpo rígido que experimenta la porción superior del elemento, dejando un brazo de palanca libre hacia la carga en la punta de:
\begin{equation}
h_{\text{lever}} = L - x_h = L - 0.05 L = 0.95 L
\end{equation}
\end{conceptbox}

% --- FIGURA 1: DIAGRAMA TIKZ DE LA COLUMNA EN VOLADIZO ---
\begin{figure}[htbp]
\centering
\begin{tikzpicture}[scale=1.15, >=Stealth]
    % Sistema de coordenadas
    \coordinate (Base) at (0,0);
    \coordinate (Hinge) at (0,0.5);
    \coordinate (Top) at (0,5.5);
    
    % Empotramiento en la base
    \draw[thick, fill=gray!30] (-1.2,-0.3) rectangle (1.2,0);
    \draw[thick] (-1.2,0) -- (1.2,0);
    \foreach \x in {-1.0,-0.7,...,1.0} {
        \draw[thick] (\x,-0.3) -- (\x-0.2,-0.55);
    }
    
    % Columna no deformada (eje de referencia)
    \draw[dashed, color=uisgray, line width=1.2pt] (0,0) -- (0,5.5);
    
    % Tramo inferior (base a rótula)
    \draw[line width=4pt, color=uisblue] (0,0) -- (0,0.5);
    
    % Resorte rotacional en 0.05 L
    \draw[color=uisgreen, fill=uisgreen!20, line width=1.5pt] (0,0.5) circle (0.28cm);
    \draw[color=uisgreen, line width=1.8pt] (-0.12,0.42) arc (200:520:0.16);
    \node[right=8pt of Hinge, font=\footnotesize\sffamily] {\textbf{\color{uisgreen}Rótula ASCE 41} ($x_h = 0.05 L$)};
    
    % Fuste deformado (curvado elásticamente + rotación de cuerpo rígido)
    \draw[line width=4pt, color=uisblue] 
        (0,0.5) to[out=86, in=-92] (1.3,5.5);
    
    % Carga lateral V en la punta
    \draw[->, color=uisrose, line width=2.5pt] (2.4,5.5) -- (1.35,5.5);
    \node[right=3pt, font=\large\bfseries\color{uisrose}] at (2.4,5.5) {$V$};
    
    % Cotas de desplazamiento en la punta
    \draw[dashed, color=uisgray] (0,5.5) -- (0,6.2);
    \draw[dashed, color=uisgray] (1.3,5.5) -- (1.3,6.2);
    \draw[<->, thick, color=uisdark] (0,6.0) -- (1.3,6.0);
    \node[above, font=\small\bfseries\color{uisdark}] at (0.65,6.0) {$\Delta = \Delta_{\text{el}} + \Delta_{\theta}$};
    
    % Cotas verticales
    \draw[<->, thick, color=uisdark] (-1.6,0) -- (-1.6,0.5);
    \node[left, font=\footnotesize] at (-1.6,0.25) {$x_h = 0.05 L$};
    
    \draw[<->, thick, color=uisdark] (-1.6,0.5) -- (-1.6,5.5);
    \node[left, font=\footnotesize] at (-1.6,3.0) {$h_{\text{lever}} = 0.95 L$};
    
    \draw[<->, thick, color=uisdark] (-2.5,0) -- (-2.5,5.5);
    \node[left, font=\small\bfseries] at (-2.5,2.75) {$L$};
    
    % Sección transversal inset
    \begin{scope}[shift={(4.0,2.2)}]
        \draw[thick, fill=gray!15] (-0.9,-1.2) rectangle (0.9,1.2);
        \draw[<->, thick] (-0.9,-1.45) -- (0.9,-1.45);
        \node[below, font=\footnotesize] at (0,-1.45) {Base $b$};
        \draw[<->, thick] (1.15,-1.2) -- (1.15,1.2);
        \node[right, font=\footnotesize] at (1.15,0) {Peralte $d$};
        % Barras de refuerzo
        \foreach \px/\py in {-0.65/-0.95, 0.65/-0.95, -0.65/0.95, 0.65/0.95, -0.65/0, 0.65/0} {
            \fill[uisdark] (\px,\py) circle (0.09);
        }
        \node[font=\footnotesize\sffamily\bfseries, color=uisdark] at (0,1.55) {Sección de Columna};
        \node[font=\scriptsize, color=uisgray] at (0,-1.85) {$I = \frac{1}{12} b d^3$};
    \end{scope}
\end{tikzpicture}
\caption{Esquema mecánico de la columna en voladizo con rótula plástica concentrada al $5\%$ de $L$ y sección rectangular de concreto armado.}
\label{fig:columna_esquema}
\end{figure}

% ==============================================================================
% SECCIÓN 2: CURVA CONSTITUTIVA ASCE 41 DE LA RÓTULA
% ==============================================================================
\section{Curva Constitutiva de la Rótula según ASCE 41}

La respuesta momento-rotación ($M_h - \theta_h$) del resorte concentrado sigue la curva envolvente generalizada (\textit{backbone curve}) definida en los capítulos de evaluación sísmica no lineal de la norma ASCE/SEI 41-17 y ASCE/SEI 41-23.

% --- FIGURA 2: CURVA ASCE 41 TIKZ ---
\begin{figure}[htbp]
\centering
\begin{tikzpicture}[scale=1.1, >=Stealth]
    % Ejes coordenados
    \draw[->, thick, color=uisdark] (-0.3,0) -- (9.5,0) node[right, font=\small\bfseries] {$\theta_h$ (rad)};
    \draw[->, thick, color=uisdark] (0,-0.3) -- (0,5.5) node[above, font=\small\bfseries] {$M_h$ (kN$\cdot$m)};
    
    % Coordenadas de los puntos ASCE 41
    \coordinate (A) at (0,0);
    \coordinate (B) at (1.5,3.2);
    \coordinate (C) at (4.5,4.2);
    \coordinate (D) at (5.0,0.85);
    \coordinate (E) at (8.0,0.85);
    \coordinate (F) at (8.8,0.1);
    
    % Umbrales IO, LS, CP
    \coordinate (IO) at (2.25, 0);
    \coordinate (LS) at (3.75, 0);
    \coordinate (CP) at (6.25, 0);
    
    % Líneas verticales de desempeño
    \draw[dashed, color=uisblue, line width=1pt] (2.25,0) -- (2.25,3.45) node[above, font=\footnotesize\bfseries, color=uisblue] {IO};
    \draw[dashed, color=uisamber, line width=1pt] (3.75,0) -- (3.75,3.95) node[above, font=\footnotesize\bfseries, color=uisamber] {LS};
    \draw[dashed, color=uisrose, line width=1pt] (6.25,0) -- (6.25,1.2) node[above, font=\footnotesize\bfseries, color=uisrose] {CP};
    
    % Trazado de la curva troncal
    \draw[line width=2pt, color=uisgreen] (A) -- (B);
    \draw[line width=2pt, color=uisamber] (B) -- (C);
    \draw[line width=2pt, color=uisrose] (C) -- (D);
    \draw[line width=2pt, color=uispurple] (D) -- (E);
    \draw[line width=1.5pt, color=uisgray, dashed] (E) -- (F);
    
    % Puntos destacados
    \foreach \p/\pos in {A/below left, B/above left, C/above, D/below left, E/above right} {
        \fill[uisdark] (\p) circle (2.8pt);
        \node[\pos, font=\bfseries\sffamily] at (\p) {\p};
    }
    
    % Líneas de referencia horizontales
    \draw[dotted, thick, color=uisgray] (0,3.2) node[left, font=\footnotesize] {$M_y$} -- (B);
    \draw[dotted, thick, color=uisgray] (0,4.2) node[left, font=\footnotesize] {$M_u = Q_y M_y$} -- (C);
    \draw[dotted, thick, color=uisgray] (0,0.85) node[left, font=\footnotesize] {$M_r = c M_y$} -- (D);
    
    % Cotas de rotación ASCE
    \draw[<->, thick, color=uisdark] (1.5,-0.6) -- (4.5,-0.6);
    \node[below, font=\footnotesize] at (3.0,-0.6) {Parámetro $a$};
    
    \draw[<->, thick, color=uisdark] (1.5,-1.3) -- (8.0,-1.3);
    \node[below, font=\footnotesize] at (4.75,-1.3) {Parámetro $b$};
    
    \draw[dotted, color=uisgray] (1.5,0) -- (1.5,-1.5);
    \draw[dotted, color=uisgray] (4.5,0) -- (4.5,-0.8);
    \draw[dotted, color=uisgray] (8.0,0) -- (8.0,-1.5);
    
    % Etiquetas de tramos
    \node[font=\scriptsize\color{uisgreen}, rotate=42] at (0.8,2.0) {Elástico ($k_{\theta,0}$)};
    \node[font=\scriptsize\color{uisamber}, rotate=14] at (3.0,3.95) {Endurecimiento};
    \node[font=\scriptsize\color{uisrose}, rotate=-78] at (4.9,2.6) {Ablandamiento ($k_t < 0$)};
    \node[font=\scriptsize\color{uispurple}] at (6.5,1.15) {Meseta Residual ($c M_y$)};
\end{tikzpicture}
\caption{Curva troncal generalizada momento-rotación de ASCE 41 con identificación de puntos canónicos A, B, C, D, E y niveles de desempeño sísmico (IO, LS, CP).}
\label{fig:curva_asce41}
\end{figure}

Los cinco puntos canónicos que gobiernan la respuesta inelástica son:
\begin{itemize}
    \item \textbf{Punto A $(0, 0)$:} Origen no deformado ni cargado.
    \item \textbf{Punto B $(\theta_y, M_y)$:} Inicio de la fluencia del refuerzo longitudinal a tracción y agrietamiento severo. La rigidez rotacional elástica inicial del resorte es:
    \begin{equation}
    k_{\theta,0} = \frac{M_y}{\theta_y}
    \end{equation}
    \item \textbf{Punto C $(\theta_u, M_u)$:} Capacidad máxima a flexión (\textit{capping point}). La rotación es $\theta_u = \theta_y + a$, donde $a$ es la capacidad de rotación plástica antes de iniciar la degradación. La sobre-resistencia se modela como $M_u = Q_y M_y$ (típicamente $Q_y \approx 1.10 - 1.20$).
    \item \textbf{Punto D $(\theta_r, M_r)$:} Inicio de la meseta residual tras una caída pronunciada de resistencia (ablandamiento post-pico con rigidez tangente negativa $k_t < 0$). La resistencia residual es $M_r = c \cdot M_y$ (para columnas con falla dúctil a flexión, $c \approx 0.20$).
    \item \textbf{Punto E $(\theta_e, M_r)$:} Límite último de deformación plástica $\theta_e = \theta_y + b$, más allá del cual ocurre pérdida total de soporte de cargas de gravedad ($M_h \to 0$, colapso seccional).
\end{itemize}

% ==============================================================================
% SECCIÓN 3: ACOPLAMIENTO CINEMÁTICO Y ECUACIONES DE EQUILIBRIO
% ==============================================================================
\section{Cinemática del Sistema en Serie y Equilibrio Estático}

El sistema columna-rótula constituye un \textbf{ensamblaje mecánico en serie} compuesto por:
\begin{enumerate}
    \item Un elemento continuo de flexión elástica con módulo de elasticidad $E$ e inercia gruesa $I = \frac{1}{12} b d^3$.
    \item Un resorte rotacional concentrado de momento $M_h$ y giro relativo $\theta_h$ en $x = x_h = 0.05 L$.
\end{enumerate}

\subsection{Flexibilidad Combinada en Serie}
Al aplicar una carga horizontal $V$ en la punta superior de la columna, el desplazamiento lateral total en la punta $\Delta$ es la suma de dos flexibilidades cinemáticas:
\begin{equation}
\Delta = \Delta_{\text{el}} + \Delta_{\theta}
\label{eq:delta_suma}
\end{equation}
donde:
\begin{itemize}
    \item $\Delta_{\text{el}}$ es la deflexión puramente elástica debida a la curvatura continua del fuste:
    \begin{equation}
    \Delta_{\text{el}} = \frac{V L^3}{3 E I}
    \label{eq:delta_el}
    \end{equation}
    \item $\Delta_{\theta}$ es la componente de desplazamiento originada por la rotación $\theta_h$ en el resorte, la cual genera una rotación de cuerpo rígido sobre la porción superior de longitud $0.95 L$:
    \begin{equation}
    \Delta_{\theta} = \theta_h \cdot (L - x_h) = \theta_h \cdot (0.95 L)
    \label{eq:delta_theta}
    \end{equation}
\end{itemize}

\subsection{Equilibrio Seccional y Diagrama de Cuerpo Libre}
Aislamos la porción superior de la columna realizando un corte transversal exactamente en el resorte rotacional ($x = x_h = 0.05 L$). Por equilibrio estático del sólido libre:
\begin{equation}
\sum M_{\text{rótula}} = 0 \implies V \cdot (0.95 L) - M_h = 0
\end{equation}
Por consiguiente, la fuerza cortante lateral $V$ y el momento flector interno en la rótula $M_h$ están vinculados de manera unívoca por:
\begin{equation}
V = \frac{M_h}{0.95 L} \iff M_h = V \cdot (0.95 L)
\label{eq:equilibrio_vh}
\end{equation}

El momento en la base empotrada ($x = 0$) tiene un brazo de palanca igual a la altura total $L$:
\begin{equation}
M_{\text{base}} = V \cdot L = \frac{M_h}{0.95 L} \cdot L = \frac{M_h}{0.95} \approx 1.0526 \cdot M_h
\label{eq:mbase}
\end{equation}
Nótese que el momento basal es siempre un $5.26\%$ superior al momento que solicita a la rótula plástica, lo cual coincide exactamente con la diferencia de brazos mecánicos ($L$ vs. $0.95 L$).

% ==============================================================================
% SECCIÓN 4: DETERMINACIÓN DEL CORTANTE BAJO CONTROL DE DESPLAZAMIENTO
% ==============================================================================
\section{Determinación del Cortante Lateral \texorpdfstring{$V$}{V} para un Desplazamiento \texorpdfstring{$\Delta$}{Delta}}

\begin{conceptbox}[Planteamiento del Problema Inverso]
En un análisis Pushover convencional, la variable independiente que se incrementa monótonamente es el \textbf{desplazamiento en la punta} $\Delta$. Por tanto, el problema fundamental consiste en responder:
\begin{center}
    \textit{Dado un desplazamiento lateral impuesto $\Delta^*$, ¿cuál es el cortante lateral equilibrado $V^*$?}
\end{center}
\end{conceptbox}

\subsection{Derivación de la Ecuación Fundamental de Compatibilidad}
Sustituyendo la relación de equilibrio \eqref{eq:equilibrio_vh} en la componente elástica \eqref{eq:delta_el}:
\begin{equation}
\Delta_{\text{el}} = \frac{\left(\frac{M_h(\theta_h)}{0.95 L}\right) L^3}{3 E I} = \frac{M_h(\theta_h) L^2}{3 \times 0.95 E I} = \frac{M_h(\theta_h) L^2}{2.85 E I}
\end{equation}
Reemplazando en la expresión cinemática total \eqref{eq:delta_suma}:
\begin{equation}
\Delta(\theta_h) = \frac{M_h(\theta_h) L^2}{2.85 E I} + 0.95 L \cdot \theta_h
\label{eq:delta_fundamental}
\end{equation}

Esta ecuación relaciona directamente el desplazamiento total impuesto $\Delta$ con la rotación interna del resorte $\theta_h$, donde $M_h(\theta_h)$ es la función constitutiva no lineal ASCE 41 detallada en la Sección~2.

\subsection{Algoritmo Numérico de Solución}
Para cualquier incremento de desplazamiento impuesto $\Delta^*$, se plantea la función residuo cinemático:
\begin{equation}
g(\theta_h) = \frac{M_h(\theta_h) L^2}{2.85 E I} + 0.95 L \cdot \theta_h - \Delta^* = 0
\label{eq:residuo_cinematico}
\end{equation}

\begin{stepbox}[Pasos del Algoritmo de Determinación de Cortante]
\begin{enumerate}
    \item \textbf{Entrada:} Desplazamiento objetivo $\Delta^*$, geometría ($L, b, d$), material ($E, f'_c$) y parámetros ASCE 41 ($M_y, \theta_y, a, Q_y, c, b$).
    \item \textbf{Búsqueda de Raíz:} Se resuelve $g(\theta_h) = 0$ para encontrar $\theta_h^*$. Dado que en secciones de columnas físicas la rigidez flexural del fuste es finita y positiva, la función $g(\theta_h)$ es estrictamente monótona creciente en el rango de interés, lo que garantiza una solución única mediante el método de bisección o Newton-Raphson:
    \begin{equation}
    \theta_h^{(k+1)} = \theta_h^{(k)} - \frac{g(\theta_h^{(k)})}{g'(\theta_h^{(k)})}
    \end{equation}
    donde la derivada tangente local es:
    \begin{equation}
    g'(\theta_h) = \frac{L^2}{2.85 E I} k_{t,\theta}(\theta_h) + 0.95 L \quad \text{con } k_{t,\theta} = \frac{dM_h}{d\theta_h}
    \end{equation}
    \item \textbf{Evaluación del Momento en la Rótula:} Con el valor convergido $\theta_h^*$, se evalúa la ley constitutiva ASCE 41:
    \begin{equation}
    M_h^* = f_{\text{ASCE}}(\theta_h^*)
    \end{equation}
    \item \textbf{Determinación del Cortante Lateral:} Por equilibrio estático riguroso:
    \begin{equation}
    V^* = \frac{M_h^*}{0.95 L}
    \end{equation}
    \item \textbf{Verificación y Descomposición:} Se recalculan las componentes individuales:
    \begin{equation}
    \Delta_{\text{el}}^* = \frac{V^* L^3}{3 E I}, \qquad \Delta_{\theta}^* = \theta_h^* \cdot (0.95 L)
    \end{equation}
    verificando que $\Delta_{\text{el}}^* + \Delta_{\theta}^* = \Delta^*$ y que el residuo de momento $|V^* \cdot 0.95 L - M_h^*| < 10^{-12}\text{ kN}\cdot\text{m}$.
\end{enumerate}
\end{stepbox}

\subsection{Rigidez Lateral Tangente del Sistema Global}
La rigidez tangente de la curva pushover ($K_t = dV / d\Delta$) se obtiene aplicando la regla de la cadena:
\begin{equation}
K_t = \frac{dV}{d\Delta} = \frac{dV}{d\theta_h} \cdot \left(\frac{d\Delta}{d\theta_h}\right)^{-1}
\end{equation}
Diferenciando las ecuaciones \eqref{eq:equilibrio_vh} y \eqref{eq:delta_fundamental}:
\begin{equation}
\frac{dV}{d\theta_h} = \frac{1}{0.95 L} k_{t,\theta}, \qquad \frac{d\Delta}{d\theta_h} = \frac{L^2}{2.85 E I} k_{t,\theta} + 0.95 L
\end{equation}
Por tanto, la rigidez lateral global tangente es:
\begin{equation}
K_t = \frac{k_{t,\theta}}{0.95 L \left( \frac{L^2}{2.85 E I} k_{t,\theta} + 0.95 L \right)} = \frac{1}{\frac{L^3}{3 E I} + \frac{(0.95 L)^2}{k_{t,\theta}}}
\label{eq:kt_global}
\end{equation}
Esta formulación en serie revela directamente que:
\begin{itemize}
    \item En el rango elástico inicial ($k_{t,\theta} = k_{\theta,0} > 0$): $K_t = K_0 > 0$.
    \item En el punto pico C ($k_{t,\theta} = 0$): el numerador se anula $\implies K_t = 0$ (fuerza máxima alcanzada).
    \item En el tramo de ablandamiento C--D ($k_{t,\theta} < 0$): $K_t < 0$, generando la pendiente descendente de la curva pushover sin divergencias numéricas gracias al control por desplazamiento.
\end{itemize}

% ==============================================================================
% SECCIÓN 5: FENÓMENO DE DESCARGA ELÁSTICA (SPRINGBACK)
% ==============================================================================
\section{Demostración del Fenómeno de Descarga Elástica (\textit{Springback})}

Uno de los conceptos físicos más sutiles y formativos en ingeniería estructural no lineal es el comportamiento del fuste elástico durante la degradación de la rótula.

\begin{alertbox}[Demostración Analítica del Enderezamiento del Fuste]
Durante la rama descendente post-pico (tramo C--D de la curva ASCE 41), el resorte rotacional experimenta ablandamiento:
\begin{equation}
k_{t,\theta} = \frac{dM_h}{d\theta_h} < 0 \implies dM_h < 0
\end{equation}
Dado que por equilibrio estricto $V = \frac{M_h}{0.95 L}$, al disminuir el momento interno resistente, la fuerza cortante lateral $V$ \textbf{debe disminuir obligatoriamente}:
\begin{equation}
dV = \frac{dM_h}{0.95 L} < 0
\end{equation}
Analicemos la derivada de la deformación elástica del fuste respecto al avance del desplazamiento total $\Delta$:
\begin{equation}
\frac{d\Delta_{\text{el}}}{d\Delta} = \frac{L^3}{3 E I} \frac{dV}{d\Delta} = \frac{L^3}{3 E I} K_t
\end{equation}
Como en el ablandamiento $K_t < 0$, se concluye formalmente que:
\begin{equation}
\frac{d\Delta_{\text{el}}}{d\Delta} < 0 \iff d\Delta_{\text{el}} < 0
\end{equation}
\textbf{Conclusión física:} ¡A pesar de que el desplazamiento total en la punta $\Delta$ continúa aumentando hacia adelante, la deflexión elástica del fuste $\Delta_{\text{el}}$ DISMINUYE! El fuste de concreto de la columna se descarga elásticamente y se endereza.
\end{alertbox}

Para compensar esta contracción elástica del fuste y al mismo tiempo acomodar el incremento de desplazamiento global positivo ($d\Delta > 0$), la rotación de la rótula $\theta_h$ debe acelerarse de forma extraordinaria:
\begin{equation}
d\Delta_{\theta} = d\Delta - d\Delta_{\text{el}} = d\Delta + |d\Delta_{\text{el}}| > d\Delta
\end{equation}
\begin{equation}
\frac{d\Delta_{\theta}}{d\Delta} = 1 - \frac{d\Delta_{\text{el}}}{d\Delta} > 1
\end{equation}
La rótula concentra toda la deformación inelástica destructiva, abriéndose con mayor velocidad para compensar la energía elástica que el fuste le transfiere al relajarse.

% ==============================================================================
% SECCIÓN 6: EJEMPLO NUMÉRICO PASO A PASO
% ==============================================================================
\section{Ejemplo Numérico Desarrollado Paso a Paso}

A continuación, se desarrolla el cálculo analítico manual de una columna típica idéntica a la calibrada en el laboratorio virtual:

\subsection{Parámetros Geométricos y Mecánicos}
\begin{itemize}
    \item Longitud de columna: $L = 3.00\text{ m}$.
    \item Sección transversal: $b = 0.35\text{ m}$ (ancho), $d = 0.45\text{ m}$ (peralte).
    \item Concreto: $f'_c = 28\text{ MPa} \implies E_c = 4700 \sqrt{28} \times 1000 = 24,870,062\text{ kN/m}^2$.
    \item Inercia gruesa:
    \begin{equation}
    I = \frac{1}{12} b d^3 = \frac{1}{12} (0.35) (0.45)^3 = 2.6578 \times 10^{-3}\text{ m}^4
    \end{equation}
    \item Rigidez a flexión: $E I = 24870062 \times 0.0026578 = 66,099.7\text{ kN}\cdot\text{m}^2$.
    \item Rigidez elástica del voladizo:
    \begin{equation}
    k_{\text{beam}} = \frac{3 E I}{L^3} = \frac{3 \times 66099.7}{27.0} = 7,344.4\text{ kN/m} = 7.3444\text{ kN/mm}
    \end{equation}
    \item Ubicación del resorte: $x_h = 0.05 L = 0.15\text{ m} \implies h_{\text{lever}} = 0.95 L = 2.85\text{ m}$.
    \item Acero y cuantía: $f_y = 420\text{ MPa}$, $\rho = 1.50\%$.
    \item Capacidad de fluencia estimada: $M_y = 379.5\text{ kN}\cdot\text{m}$.
    \item Rotación de fluencia del resorte: $\theta_y = 0.0020\text{ rad} \implies k_{\theta,0} = \frac{379.5}{0.0020} = 189,750\text{ kN}\cdot\text{m/rad}$.
    \item Parámetros ASCE 41:
    \begin{itemize}
        \item $a = 0.020\text{ rad}$, $Q_y = 1.15 \implies M_u = 1.15 \times 379.5 = 436.4\text{ kN}\cdot\text{m}$, $\theta_u = \theta_y + a = 0.0220\text{ rad}$.
        \item $c = 0.20 \implies M_r = 0.20 \times 379.5 = 75.9\text{ kN}\cdot\text{m}$, $\theta_r = 0.0250\text{ rad}$.
        \item $b = 0.040\text{ rad} \implies \theta_e = \theta_y + b = 0.0420\text{ rad}$.
    \end{itemize}
\end{itemize}

Constante geométrica de compatibilidad:
\begin{equation}
C_{\text{el}} = \frac{L^2}{2.85 E I} = \frac{(3.0)^2}{2.85 \times 66099.7} = 4.7774 \times 10^{-5}\text{ m/(kN}\cdot\text{m)}
\end{equation}
La ecuación de desplazamiento \eqref{eq:delta_fundamental} se reduce a:
\begin{equation}
\Delta(\theta_h) = 4.7774 \times 10^{-5} M_h(\theta_h) + 2.85 \cdot \theta_h \quad \text{[m]}
\end{equation}

% --- PASO 1 ---
\subsection{Paso 1: Rango Elástico Lineal (\texorpdfstring{$\Delta = 10.0\text{ mm} = 0.010\text{ m}$}{Delta = 10.0 mm})}

\begin{examplebox}[Cálculo Detallado del Paso 1]
Como $\Delta$ es pequeño, asumimos comportamiento elástico de la rótula: $M_h = k_{\theta,0} \theta_h = 189750 \theta_h$.
Sustituyendo en la ecuación cinemática:
\begin{align}
0.010 &= 4.7774 \times 10^{-5} (189750 \theta_h) + 2.85 \theta_h \nonumber\\
0.010 &= 9.0651 \theta_h + 2.85 \theta_h = 11.9151 \theta_h
\end{align}
Despejamos el giro en la rótula:
\begin{equation}
\theta_h = \frac{0.010}{11.9151} = 8.3927 \times 10^{-4}\text{ rad} = 0.8393\text{ mrad}
\end{equation}
Verificación de rango: $\theta_h = 0.000839 < \theta_y = 0.0020\text{ rad} \implies$ \textbf{¡Correcto, rango elástico A--B!}

Calculamos el momento en la rótula:
\begin{equation}
M_h = 189750 \times 8.3927 \times 10^{-4} = 159.25\text{ kN}\cdot\text{m}
\end{equation}
Calculamos la fuerza cortante lateral por equilibrio:
\begin{equation}
V = \frac{M_h}{0.95 L} = \frac{159.25}{2.85} = \mathbf{55.88\text{ kN}}
\end{equation}
Momento en la base:
\begin{equation}
M_{\text{base}} = V \cdot L = 55.88 \times 3.0 = \mathbf{167.64\text{ kN}\cdot\text{m}} = \frac{159.25}{0.95}
\end{equation}
Descomposición de desplazamientos:
\begin{align}
\Delta_{\text{el}} &= \frac{V L^3}{3 E I} = \frac{55.88 \times 27}{3 \times 66099.7} = 7.61 \times 10^{-3}\text{ m} = \mathbf{7.61\text{ mm}}\\
\Delta_{\theta} &= \theta_h \times 2.85 = 8.3927 \times 10^{-4} \times 2.85 = 2.39 \times 10^{-3}\text{ m} = \mathbf{2.39\text{ mm}}
\end{align}
Verificación de compatibilidad y equilibrio:
\begin{equation}
\Delta_{\text{el}} + \Delta_{\theta} = 7.61 + 2.39 = \mathbf{10.00\text{ mm}} \quad (\text{Error} = 0.000\text{ mm})
\end{equation}
\begin{equation}
R_M = V \cdot (0.95 L) - M_h = 55.88 \times 2.85 - 159.25 = \mathbf{0.00\text{ kN}\cdot\text{m}} \quad \checkmark
\end{equation}
\end{examplebox}

% --- PASO 2 ---
\subsection{Paso 2: Fluencia e Inicio de Plastificación (\texorpdfstring{$\theta_h = \theta_y = 0.0020\text{ rad}$}{theta = 0.0020 rad})}

\begin{examplebox}[Cálculo Detallado del Paso 2 (Punto B de ASCE 41)]
En el umbral de plastificación, el momento alcanza la fluencia seccional:
\begin{equation}
M_h = M_y = 379.50\text{ kN}\cdot\text{m}
\end{equation}
El cortante lateral en fluencia es:
\begin{equation}
V_y = \frac{M_y}{0.95 L} = \frac{379.50}{2.85} = \mathbf{133.16\text{ kN}}
\end{equation}
Momento basal en fluencia:
\begin{equation}
M_{\text{base},y} = 133.16 \times 3.0 = \mathbf{399.47\text{ kN}\cdot\text{m}}
\end{equation}
Desplazamiento global de fluencia en la punta:
\begin{align}
\Delta_y &= \frac{V_y L^3}{3 E I} + \theta_y \times 2.85 \nonumber\\
&= \frac{133.16 \times 27}{3 \times 66099.7} + 0.0020 \times 2.85 \nonumber\\
&= 0.01813\text{ m} + 0.00570\text{ m} = 0.02383\text{ m} = \mathbf{23.83\text{ mm}}
\end{align}
Rigidez lateral inicial elástica global:
\begin{equation}
K_0 = \frac{V_y}{\Delta_y} = \frac{133.16\text{ kN}}{23.83\text{ mm}} = \mathbf{5.588\text{ kN/mm}} = 5,588\text{ kN/m}
\end{equation}
\end{examplebox}

% --- PASO 3 ---
\subsection{Paso 3: Endurecimiento Plástico (\texorpdfstring{$\Delta = 40.0\text{ mm} = 0.040\text{ m}$}{Delta = 40.0 mm})}

\begin{examplebox}[Cálculo Detallado del Paso 3 (Tramo B--C)]
Como $\Delta = 40.0\text{ mm} > \Delta_y = 23.83\text{ mm}$, la rótula ha plastificado y se encuentra en la rama de endurecimiento plástico (tramo B--C).
La ley de endurecimiento entre $B$ ($0.002, 379.5$) y $C$ ($0.022, 436.4$) tiene una pendiente tangente:
\begin{equation}
k_{t,\theta} = \frac{M_u - M_y}{\theta_u - \theta_y} = \frac{436.425 - 379.50}{0.020} = \frac{56.925}{0.020} = 2,846.25\text{ kN}\cdot\text{m/rad}
\end{equation}
La ecuación constitutiva en este tramo es:
\begin{equation}
M_h(\theta_h) = 379.50 + 2846.25 (\theta_h - 0.0020) = 373.8075 + 2846.25 \theta_h
\end{equation}
Sustituimos en la ecuación de compatibilidad cinemática:
\begin{align}
0.040 &= 4.7774 \times 10^{-5} (373.8075 + 2846.25 \theta_h) + 2.85 \theta_h \nonumber\\
0.040 &= 0.017858 + 0.13598 \theta_h + 2.85 \theta_h \nonumber\\
0.022142 &= 2.98598 \theta_h \implies \theta_h = \mathbf{7.4153 \times 10^{-3}\text{ rad}} = \mathbf{7.42\text{ mrad}}
\end{align}
Momento resistente en la rótula:
\begin{equation}
M_h = 379.50 + 2846.25 (0.0074153 - 0.0020) = \mathbf{394.91\text{ kN}\cdot\text{m}}
\end{equation}
Cortante lateral equilibrado:
\begin{equation}
V = \frac{M_h}{2.85} = \frac{394.91}{2.85} = \mathbf{138.56\text{ kN}}
\end{equation}
Descomposición de desplazamientos:
\begin{align}
\Delta_{\text{el}} &= \frac{138.56 \times 27}{3 \times 66099.7} = 0.01887\text{ m} = \mathbf{18.87\text{ mm}}\\
\Delta_{\theta} &= 0.0074153 \times 2.85 = 0.02113\text{ m} = \mathbf{21.13\text{ mm}}\\
\Delta_{\text{total}} &= 18.87 + 21.13 = \mathbf{40.00\text{ mm}} \quad (\text{Equilibrio y compatibilidad exactos})
\end{align}
Rigidez tangente global en endurecimiento:
\begin{align}
K_t &= \frac{1}{\frac{1}{7.3444} + \frac{(2.85)^2}{2.84625}} = \frac{1}{0.13616 + 2.85375} = \mathbf{0.334\text{ kN/mm}} \nonumber\\
    &\quad (\text{Degradación severa de rigidez: } K_t \approx 0.06 K_0)
\end{align}
\end{examplebox}

% --- PASO 4 ---
\subsection{Paso 4: Ablandamiento y Descarga Elástica (\texorpdfstring{$\Delta = 100.0\text{ mm} = 0.100\text{ m}$}{Delta = 100.0 mm})}

\begin{examplebox}[Cálculo Detallado del Paso 4 (Meseta Residual y Comprobación de Springback)]
A $\Delta = 100\text{ mm}$, la rótula ha superado el tramo de ablandamiento post-pico C--D y ha ingresado en la meseta residual D--E:
\begin{equation}
M_h = M_r = c \cdot M_y = 0.20 \times 379.50 = \mathbf{75.90\text{ kN}\cdot\text{m}}
\end{equation}
El cortante lateral se reduce drásticamente por pérdida de capacidad en la rótula:
\begin{equation}
V = \frac{M_r}{2.85} = \frac{75.90}{2.85} = \mathbf{26.63\text{ kN}}
\end{equation}
Calculamos la deformación elástica del fuste para este cortante residual:
\begin{equation}
\Delta_{\text{el}} = \frac{V L^3}{3 E I} = \frac{26.63 \times 27}{3 \times 66099.7} = 3.63 \times 10^{-3}\text{ m} = \mathbf{3.63\text{ mm}}
\end{equation}
\textbf{¡Demostración de Springback!:}\\
En el Paso 3 ($\Delta = 40\text{ mm}$), la deformación elástica del fuste era $\Delta_{\text{el}} = 18.87\text{ mm}$. En el pico C llegó a ser $\approx 19.8\text{ mm}$. Ahora, a $\Delta = 100\text{ mm}$, $\Delta_{\text{el}}$ ha caído a solo \textbf{3.63 mm}. ¡El fuste ha recuperado más de $16\text{ mm}$ de deflexión elástica enderezándose!

El aporte de rotación en la rótula plástica debe absorber todo el resto del desplazamiento:
\begin{equation}
\Delta_{\theta} = \Delta - \Delta_{\text{el}} = 100.0 - 3.63 = \mathbf{96.37\text{ mm}} = 0.09637\text{ m}
\end{equation}
El giro en la rótula se dispara a:
\begin{equation}
\theta_h = \frac{\Delta_{\theta}}{2.85} = \frac{0.09637}{2.85} = \mathbf{0.03381\text{ rad}} = \mathbf{33.81\text{ mrad}}
\end{equation}
Verificación de rango: $\theta_r = 0.0250 < \theta_h = 0.03381 < \theta_e = 0.0420\text{ rad} \implies$ \textbf{¡Meseta residual D--E confirmada!}

Residuo numérico exacto:
\begin{equation}
R_M = 26.63 \times 2.85 - 75.90 = 75.90 - 75.90 = \mathbf{0.000\text{ kN}\cdot\text{m}} \quad \checkmark
\end{equation}
\end{examplebox}

% ==============================================================================
% SECCIÓN 7: TABLA COMPARATIVA DE RESULTADOS
% ==============================================================================
\section{Tabla Comparativa de Resultados en las Fases de Deformación}

La Tabla~\ref{tab:resumen_pushover} resume la evolución de las variables mecánicas a lo largo del proceso de empuje Pushover con control por desplazamiento.

\begin{table}[htbp]
\centering
\small
\renewcommand{\arraystretch}{1.25}
\resizebox{\linewidth}{!}{%
\begin{tabular}{lcccccccc}
\toprule
\textbf{Fase / Punto} & $\bm{\Delta}$ [mm] & \textbf{Estado ASCE} & $\bm{\theta_h}$ [mrad] & $\bm{M_h}$ [kN$\cdot$m] & $\bm{V}$ [kN] & $\bm{M_{\text{base}}}$ [kN$\cdot$m] & $\bm{\Delta_{\text{el}}}$ [mm] & $\bm{\Delta_{\theta}}$ [mm] \\
\midrule
1. Reposo & 0.0 & Elástico A & 0.00 & 0.0 & 0.0 & 0.0 & 0.00 & 0.00 \\
2. Elástico & 10.0 & Elástico A--B & 0.84 & 159.3 & 55.9 & 167.6 & 7.61 & 2.39 \\
3. Fluencia (B) & 23.8 & Umbral Fluencia & 2.00 & 379.5 & 133.2 & 399.5 & 18.13 & 5.70 \\
4. Endurecimiento & 40.0 & Plástico B--C & 7.42 & 394.9 & 138.6 & 415.7 & 18.87 & 21.13 \\
5. Capacidad Pico (C) & 82.5 & Capping Point & 22.00 & 436.4 & 153.1 & 459.4 & 19.82 & 62.70 \\
6. Ablandamiento & 86.0 & Degradación C--D & 23.50 & 256.2 & 89.9 & 269.7 & 12.22 & 66.98 \\
7. Residual (D--E) & 100.0 & Meseta Residual & 33.81 & 75.9 & 26.6 & 79.9 & \textbf{3.63} & \textbf{96.37} \\
8. Límite Último (E) & 123.3 & Colapso Inminente & 42.00 & 75.9 & 26.6 & 79.9 & 3.63 & 119.70 \\
\bottomrule
\end{tabular}%
}
\caption{Evolución paso a paso de fuerzas, momentos y deformaciones. Nótese la marcada reducción de $\Delta_{\text{el}}$ de 19.82 mm a 3.63 mm durante el ablandamiento post-pico (descarga elástica o \textit{springback}).}
\label{tab:resumen_pushover}
\end{table}


% ==============================================================================
% SECCIÓN 8: CONCLUSIONES Y RECOMENDACIONES
% ==============================================================================
\section{Conclusiones Pedagógicas y de Diseño}

\begin{enumerate}
    \item \textbf{Equivalencia Cinemática de la Plasticidad Concentrada:} Ubicar el resorte rotacional no lineal al $5\%$ de la altura ($x_h = 0.05 L$) representa rigurosamente el centroide del bloque de curvatura inelástica de longitud $L_p \approx 0.10 L$, desacoplando la flexibilidad de flexión del fuste de la rotación concentrada de la rótula.
    \item \textbf{Resolución del Cortante bajo Control por Desplazamiento:} La fuerza lateral cortante $V$ no se puede imponer libremente cuando la estructura incursiona en el rango post-pico. En su lugar, el desplazamiento impuesto $\Delta$ actúa como variable de control cinemático, determinando unívocamente la rotación $\theta_h$ por compatibilidad en serie y fijando el cortante $V = M_h(\theta_h) / (0.95 L)$ a través del equilibrio estricto de momentos seccionales.
    \item \textbf{Naturaleza Física del \textit{Springback}:} Durante la etapa de degradación post-capping ($k_{t,\theta} < 0$), la caída de resistencia interna en la rótula provoca que el cortante $V$ disminuya. Esta disminución de fuerza descarga elásticamente al fuste de la columna, haciendo que su curvatura continua disminuya y obligando a la rótula a abrirse aceleradamente. Ignorar este acoplamiento conduce a estimaciones erróneas de la demanda de ductilidad local.
    \item \textbf{Criterios de Desempeño ASCE 41:} Los límites IO, LS y CP brindan umbrales objetivos de daño. El modelo demuestra que entre la fluencia (B) y el límite Life Safety (LS), la estructura conserva estabilidad global con endurecimiento positivo, mientras que superar el punto C sitúa al elemento en una zona de alta vulnerabilidad donde la resistencia residual depende exclusivamente de la integridad del confinamiento del núcleo.
\end{enumerate}

% ==============================================================================
% REFERENCIAS BIBLIOGRÁFICAS
% ==============================================================================
\begin{thebibliography}{99}
    \bibitem{asce41} American Society of Civil Engineers (ASCE). (2017). \textit{Seismic Evaluation and Retrofit of Existing Buildings} (ASCE/SEI Standard 41-17). Reston, VA: ASCE.
    \bibitem{fema356} Federal Emergency Management Agency (FEMA). (2000). \textit{Prestandard and Commentary for the Seismic Rehabilitation of Buildings} (FEMA 356). Washington, DC: FEMA.
    \bibitem{chopra} Chopra, A. K. (2020). \textit{Dynamics of Structures: Theory and Applications to Earthquake Engineering} (5th ed.). Hoboken, NJ: Pearson.
    \bibitem{priestley} Paulay, T., \& Priestley, M. J. N. (1992). \textit{Seismic Design of Reinforced Concrete and Masonry Buildings}. New York: John Wiley \& Sons.
    \bibitem{powell} Powell, G. H. (2010). \textit{Nonlinear Analysis of Concrete and Steel Structures: An Introduction to the Mechanics and Analysis of Inelastic Structures}. Berkeley, CA: Computers and Structures, Inc.
\end{thebibliography}

\end{document}
